
\section{Conductor ideals in valuation domains}

The following statement includes valuation groups of arbitrary rank and
does not require a coefficient field in the base ring.

\begin{theorem}\label{thm:valuation-conductor-arbitrary-coefficients}
Let $V$ be a valuation domain with fraction field $K$ and infinite
residue field $l$. Suppose that $R\subset V$ contains a nonzero proper
$V$-ideal $J$ and that $R/J$ is a finite field. For every inclusion of
commutative rings $R\subset S$, the following conditions are equivalent:
\[
\begin{gathered}
 \text{the standard coaddition and comultiplication tensors exist uniquely},\\
 \theta_2\text{ is injective},\qquad (R:_S J)=R,\\
 \operatorname{Ann}_{S/R}(J)=0,\qquad \Omega(R,S)=R[X].
\end{gathered}
\]
In these cases all multivariable tensor maps are isomorphisms, and the
compatible plethory is the ordinary polynomial plethory.
\end{theorem}

We first establish the valuation-ideal calculation used in the proof.
It is a form of the classical Clifford regularity of ideals in a
valuation domain; see Halter-Koch~\cite[Section~5, Lemmas~5.1--5.3 and
Theorem~5.4(1)]{HalterKoch2009}.

\begin{lemma}\label{lem:valuation-ideal-division}
For every nonzero ideal $J$ of a valuation domain and every $N\ge1$,
\[
                         J^N(J:_KJ^N)=J.
\]
\end{lemma}
\begin{proof}
Put $U=(J:J)$, a valuation overring, and let $\mathfrak p$ be its
maximal ideal. If $J=aU$, then $J=J^2a^{-1}U$. Otherwise
$J\mathfrak p=J$: an element $a\in J\setminus J\mathfrak p$ would
generate $J$, because $b\notin aU$ implies $a/b\in\mathfrak p$.
In the nonprincipal case, $J(U:J)=\mathfrak p$. This product is proper,
since equality with $U$ would make $J$ invertible and therefore
principal over the local valuation ring. Conversely, for
$z\in\mathfrak p$, equality $zJ=J$ would imply $z^{-1}\in U$.
Choose $a\in J\setminus zJ$. Comparability gives $zJ\subset aU$,
so $z/a\in(U:J)$ and $z\in J(U:J)$.
Thus in either case $J=J^2A$ for a fractional ideal $A$. Iterating
gives $J=J^NA^{N-1}$, and $A^{N-1}\subset(J:J^N)$. This proves one
inclusion; the other is the definition of the ideal quotient.
\end{proof}

\begin{proof}[Proof of Theorem~\ref{thm:valuation-conductor-arbitrary-coefficients}]
The ring $R$ is local with maximal ideal $J$. Indeed $R/J$ injects
into $V/\mathfrak m_V$, and the inverse of $1+j$, for $j\in J$, lies
in $1+J$. Write $q=|R/J|$.

We first dispose of the idempotent case. If $J^2=J$, then for every
$N$ the ideal $J$ is generated as an $R$-ideal by the $N$th powers
of its elements. To see this, write $0\ne r\in J=J^{N+1}$ as a
finite sum of products and choose a product of value at most $v(r)$.
For its factor $y$ of smallest value, $r/y^N$ belongs to $J$.
For such a local ring the Vandermonde ideal $D_d(R)$ is $R$ when
$d<q$ and is $J$ when $d\ge q$. The first assertion uses distinct
residue classes. For the second, all determinants belong to $J$;
evaluating the determinant polynomial on $0,1,a,a^2,\ldots$ gives
$a^{d(d-1)(d+1)/6}$ times a unit for $a\in J$, and the generation
by powers gives the reverse inclusion. Consequently the effective
coefficient ring is $C=(R:_S J)$. If $R\to C$ is epic, then $JC$ is
an ideal of $R$ containing $J$. It cannot equal $R$: an equation
$1=\sum a_ic_i$ would make $J$ finitely generated, and $J=J^2$ would
contradict Nakayama's lemma. Hence $JC=J$. Reducing the epimorphism
modulo $J$ gives the injective field epimorphism $R/J\to C/J$, so
$C=R$. The coefficient-ring obstruction already proved then gives
the theorem in this case.

Assume now $J^2\ne J$ and fix $0\ne c\in J\setminus J^2$. Set
$W=V[1/c]=R[1/c]$. For every $0\ne y\in J$ one has
$2v(y)>v(c)$; comparison of the factors in a product therefore gives
\begin{equation}\label{eq:conductor-power-control}
                        J^{2h+1}\subset c^hJ\quad(h\ge0).
\end{equation}
The ideals $D_d(R)$ are cofinal with the powers of $J$. Many nodes
in the $q$ residue classes force arbitrarily many differences into
$J$, proving one direction. For the other, choose $d+1$ elements
$u_i\in V$ with distinct residues. The determinant on $cu_i$ is
$\delta=c^hu$ with $h=d(d+1)/2$ and $u\in V^\times$; then
$J^{2h+1}\subset c^hJ=\delta J\subset\delta R$.
Thus the effective coefficient ring is
\[
                            C=\bigcup_{n\ge1}(R:_S J^n).
\]
Iterating colons shows that $(R:_S J)=R$ implies $C=R$, which proves
all positive conclusions. Suppose for the rest of the proof that
$\theta_2$ is injective. The coefficient-ring obstruction makes
$R\to C$ an epimorphism.

Rationalization has image $L\subset W$ and kernel $T$ annihilated
elementwise by powers of $J$. If every $b\in L$ has
$v(b)\ge-v(c)$, then $J^3L\subset J$ is a common ideal containing
$J^3$. The injective epimorphism
$R/J^3L\to L/J^3L$ is an isomorphism: reduction by the nilpotent
maximal ideal gives an epimorphism from a field, and finite iteration
of that ideal proves surjectivity. Hence $L=R$. Otherwise choose
$b\in L$ with $v(b)<-v(c)$. Since $b^{-1}V\subset J$, one obtains
$V\subset L$ and $c^{-1}=b(cb)^{-1}\in L$, whence $L=W$.
The case $L=R$ retracts the epimorphism $R\to C$ and gives $C=R$.

In the remaining case lift $c^{-1}$ to $a\in C$ and put $z=ca$.
Some $c^n$ annihilates $z-1$, so $z^n(z-1)=0$. Therefore $e=z^n$
is idempotent and $e=z^N$ for all $N\ge n$. It annihilates $T$,
and rationalization identifies $eC$ with $W$. The element $1-e$ is
killed by some $J^M$. The same nilpotent-ideal argument on the other
factor gives
\begin{equation}\label{eq:valuation-coefficient-splitting}
                 C=W\times R/I,\qquad J^M\subset I\subset R,
\end{equation}
with the diagonal inclusion of $R$; $I=R$ is allowed.

Choose representatives $B\subset R$ of $R/J$ and put
$P(X)=\prod_{b\in B}(X-b)$. Differences of representatives are
units, $P(R)\subset J$, and $P'(b)$ is a unit for every $b\in B$.
The monic basis $X^iP^j$, $0\le i<q$, gives the exact decomposition
\begin{equation}\label{eq:valuation-conductor-coefficients}
 \Omega(R,K)=\bigoplus_{i<q}RX^i\ \oplus\!
 \bigoplus_{\substack{i<q\\j\ge1}}(J:_KJ^j)X^iP^j.
\end{equation}
Here is the necessity argument, including the boundary of any
valuation cut. Subtract interpolation on $B$, leaving only $j\ge1$.
For fixed $0\ne y\in J$, divide the substitution $X=b+yZ$ by a
coefficient of minimum value among the $a_{ij}y^j$. Reduction gives
\[
 \sum_{j\ge1}\left(\sum_{i<q}\overline{a_{ij}y^j/\mu}\,\bar b^i\right)
                         (\overline{P'(b)}Z)^j.
\]
For some $b$ this is a nonzero polynomial, by interpolation on the
finite field, and it has zero constant term. A negative minimum
contradicts $V$-valuedness. A zero minimum would give a nonconstant
polynomial from the infinite field $l$ into a finite set. A positive
minimum outside $v(J)$ gives a value in $R\cap\mathfrak m_V$ outside
$J$. Thus $a_{ij}y^j\in J$. Products of $j$ elements of $J$ are
divisible by the $j$th power of their least-valued factor, proving
$a_{ij}J^j\subset J$. Sufficiency is immediate. All the coefficient
ideals displayed above lie in $W$, since $c^ja_{ij}\in J$.

For $F$ vanishing on $B$, define the ideal
\[
              \nu(F)=\sum_{i<q,j\ge1}a_{ij}J^j.
\]
If $F(b+Y)=\sum_{h\ge1}g_{bh}Y^h$, the same interpolation calculation
at each fixed $y\in J$ proves
\begin{equation}\label{eq:valuation-ideal-norm}
                         \nu(F)=\sum_{b\in B,h\ge1}g_{bh}J^h.
\end{equation}
Indeed the two sides have equal least coefficient values after
replacing every $J^h$ by $y^hV$; the union over $y$ computes the
powers of a valuation ideal. No numerical infimum is used.

If $I=R$ in \eqref{eq:valuation-coefficient-splitting}, the coefficient
module of $P$ in \eqref{eq:valuation-conductor-coefficients} is
$U=(J:J)$, with $JU=J$. Choose $\beta\in V$ whose residue is not in
$R/J$. The tensor $\beta P\otimes P-P\otimes\beta P$ has zero
polynomial image. Projection to $U\otimes_RU$ and reduction modulo
$J$ give $\bar\beta\otimes1-1\otimes\bar\beta\ne0$ over the field
$R/J$. Hence this case contradicts injectivity.

Suppose $I\subsetneq R$. Put
\[
\begin{aligned}
 H&=\{F\in W[X]:F(R)\subset I\},& E_0&=R[X]+H,\\
 H_0&=\{F\in H:F|_B=0\},& K_0&=\sum_{F\in H_0}\nu(F).
\end{aligned}
\]
The first projection of $\Omega(R,C)$ is $E_0$, by lifting its
second component to $R[X]$. Interpolation on $B$ maps $H$ to $H_0$,
and the polynomials $aP$, $a\in J^{M-1}$, show
$J^M\subset K_0\subset J$.

We claim $H\not\subset JE_0$. Otherwise, for $F\in H_0$ and
$N\ge1$, substitute
$\psi=b+\lambda P^N$, where $b\in B$ and $\lambda\in(J:J^N)$.
This preserves $R$, and $F\circ\psi\in H_0$ has only $P$-degrees
at least $N$. Lemma~\ref{lem:valuation-ideal-division}, pure-power
generation of powers of valuation ideals, and
\eqref{eq:valuation-ideal-norm} give
\[
 \sum_{b,\lambda}\nu(F\circ\psi)
   =\sum_{b,h}g_{bh}(J:J^N)^hJ^{Nh}
   =\sum_{b,h}g_{bh}J^h=\nu(F).
\]
Write $F\circ\psi=A+\sum r_tF_t$, with $A\in JR[X]$, $r_t\in J$,
and $F_t\in H$. Subtract their interpolants to arrange $F_t\in H_0$.
Projection to $P$-degrees at least $N$ gives
$K_0\subset J^{N+1}+JK_0$. Taking $N=M$ yields $K_0=JK_0$.
Iteration would imply $J^M\subset K_0\subset J^{M+1}$.
But $J^M=J^{M+1}$ forces $J=J^2$: express $r^M\in J^{2M}$ as a
finite sum of products and choose a least-valued factor $y$ of a
product of value at most $Mv(r)$; then $2v(y)\le v(r)$ and
$r\in y^2V$. This contradicts our choice of $c$, proving the claim.

Choose $F\in H\setminus JE_0$. The two integer-valued polynomials
$f=(F,0)$ and $g=(0,P^M\bmod I)$ survive modulo $J\Omega(R,C)$:
use the first projection for $f$ and the monic coefficient of the
second projection for $g$. Their tensor remains nonzero over the
residue field, while its polynomial image is zero because their
coefficient supports are disjoint. This is the final contradiction.
Hence injectivity forces $C=R$. Iterated colons and the universal
coefficient description finish every asserted equivalence.
\end{proof}

\begin{corollary}\label{cor:valuation-conductor-specializations}
The theorem applies, with arbitrary coefficient torsion, to
$R=\mathbb F_q+t^al[[t]]$ for any infinite extension
$l/\mathbb F_q$ and $a\ge1$; to $R=\mathbb F_q+xV$ for an arbitrary
valuation domain of characteristic $p$, infinite residue field, and
$0\ne x\in\mathfrak m_V$; and to $R=\mathbb F_q+J$ for every
nonzero proper ideal in such a valuation domain. In particular no
attained or rational cut is required for a real-valued valuation,
and no rank restriction is required in the general case.
\end{corollary}

